\chapter{Programmablaufprotokolle}

\section{Steuerung}

Die Spielfigur wird durch das anklicken der gelben Pfeile an den Rändern des Spielfensters bzw. der Türen durch das Schulhaus bewegt. Wird die Spielfigur selbst angeklickt, öffnet sich deren Inventar
und bereits gesammelte Items können eingesehen werden. Das Inventar kann mit einem Klick auf sein rotes \glqq{}X\grqq{} wieder geschlossen werden.

Wird die Taste \texttt{P} gedrückt, öffnet sich ein Fenster, in welchem die Farben der Haare, der
Haut, des Hoodies und der Hose der Spielfigur angepasst werden kann. Diese Farben können durch
selbiges Fenster ebenfalls als Datei gespeichert bzw. wieder geladen werden.

\subsection{Minigames}

\subsubsection{Informatik}

\subsubsection{Mathematik}

\subsubsection{Chemie}

\subsubsection{Physik}

\subsubsection{Biologie}


\subsection{Kampf}


\section{Tests}

\subsection{\texttt{Unit Testing}}

Mithilfe des \texttt{JUnit}-Frameworks wurde eine Reihe von Tests implementiert,
welche die Funktionalität des Szenenstapels, also dessen Operationen, sowie das
Szenen-Verdeckungs-Einstellungs-System testet. Die Tests werden bei jedem Push
auf die GitHub-Repository ausgeführt, es ist also gleich ersichtlich, wenn eine
Änderung unbeabsichtigte Nebeneffekte hatte.

\subsection{Testen während der Entwicklung}

Bei der Implementierung von neuen Funktionen wurden diese stets umgehend soweit
getestet und ggf. geändert, dass sie allgemein dem Erwartungsbild entsprachen.
Andernfalls wurden einige Fehler im späteren Verlauf der Entwicklung beim Spielen
des Spiels aufgedeckt, gemeldet und, soweit möglich, behoben.
